\begin{figure}[ht]
\centering
\begin{tikzpicture}[>=Stealth, font=\small, node distance=1.5cm]
  \node[draw, rounded corners, fill=blue!8, minimum width=4.2cm, minimum height=0.9cm] (a) {习题 1--2：亏格 $1$ 与 $J_0(11)$};
  \node[draw, rounded corners, fill=green!8, below=of a, minimum width=4.2cm, minimum height=0.9cm] (b) {习题 3：Hecke 对应与 Rosati};
  \node[draw, rounded corners, fill=orange!10, below=of b, minimum width=4.2cm, minimum height=0.9cm] (c) {习题 4：特征形式因子分解};
  \node[draw, rounded corners, fill=purple!10, below=of c, minimum width=4.2cm, minimum height=0.9cm] (d) {习题 5：二维 Jacobian $J_0(37)$};

  \draw[->, thick] (a) -- (b);
  \draw[->, thick] (b) -- (c);
  \draw[->, thick] (c) -- (d);
\end{tikzpicture}
\caption{本章习题路线图：先巩固“亏格 $1$ 曲线就是自己的 Jacobian”，再练习 Hecke 对应与极化，最后把不同特征形式对应的 Abel 因子从 $J_0(N)$ 中拆出来。}
\label{fig:ch06-exercises-roadmap}
\end{figure}
